\section{Final Design Report}
%=============================================================================

\begin{faqbox}[Q: How much of the final report can come from earlier sprint submissions?]
Much of it \textit{should} come from earlier sprints; that is by design. However, ``copy and paste from Sprint~3'' is not sufficient. You are expected to \textbf{revise and update} every section to reflect the final state of your project and address any issues identified in earlier sprint feedback. Treat this as a polished, standalone document that tells the complete story.
\end{faqbox}

\begin{faqbox}[Q: What goes in the appendix vs.\ the main body?]
The main body should contain your narrative, key results, and high-level design information. The appendix is for detailed supporting material: full CAD drawings, complete BOM, raw test data, detailed calculations, wiring diagrams, and software documentation. A good rule of thumb: if removing it would break the narrative flow of the main body, it belongs in the main body. If it is supporting evidence or reference material, it belongs in the appendix.
\end{faqbox}

\begin{protip}
\textbf{From Evidence to Narrative.} The final report is not a compilation of sprint deliverables stapled together. It is a \textit{story} with a structure:

\begin{enumerate}[leftmargin=*, itemsep=1pt, topsep=0pt]
    \item \textbf{Introduction} sets the scene: the problem, the stakeholders, why it matters.
    \item \textbf{Design Evolution} is the rising action: concept selection, key decisions, what changed and why.
    \item \textbf{Results} is the climax: VCRM verdicts, test data, system performance vs.\ requirements.
    \item \textbf{Lessons Learned} is the resolution: what worked, what did not, what comes next.
\end{enumerate}

Each sprint's deliverables become \textit{supporting evidence} in the appendix, not the main text. In the body, summarize: ``The power subsystem was redesigned between CDR and FDR to address a motor inrush brown-out (see Appendix~C for the full change log). The revised design eliminated the 3.3\,V rail sag measured during Sprint~4 testing.'' Do not paste in the raw Sprint~4 report. A reader should be able to understand the full project story from the main body alone, then consult the appendices for proof.
\end{protip}

\begin{tipbox}[Map the Narrative to the Required Report Sections]
The Design Project spec defines 7 numbered final report sections. Make sure the narrative arc above covers all of them; it is easy to miss the standalone sections that do not fit neatly into the story:
\begin{enumerate}[nosep,leftmargin=1.5em]
\item \textbf{Introduction} (problem, stakeholders, background)
\item \textbf{Design Evolution} (concept selection, key decisions, as-designed vs.\ as-built)
\item \textbf{Hazards Discussion} (FMEA update, highest-RPN items, mitigations applied; often omitted by accident)
\item \textbf{Future Work} (v2.0 improvements, DfD, O\&M, lifecycle considerations; also frequently omitted)
\item \textbf{Results} (VCRM verdicts, test data, system performance vs.\ requirements)
\item \textbf{Budget} (as-designed vs.\ as-built cost, variances explained)
\item \textbf{Lessons Learned} (specific, actionable reflections)
\end{enumerate}
Sections~3 (Hazards Discussion) and 4 (Future Work) are the most commonly omitted because they stand alone rather than emerging naturally from the narrative. Build them into your outline from the start.
\end{tipbox}

\begin{faqbox}[Q: What should go in the ``Future Work'' and ``Hazards Discussion'' sections?]
For \textbf{Future Work}: think in terms of lifecycle and sustainability. The Master Reference Document includes a formal Design for Disassembly requirement: ``The system shall be fully disassembled into its base material categories using only standard hand tools in under 5 minutes.'' This is testable at FDR by demonstration. If your design uses permanent adhesives where fasteners would work, document why and propose a DfD-compliant redesign for the next version. What refinements would bring the prototype to a production-ready product? Consider serviceability (can a technician access and replace the most failure-prone components?), maintenance planning (what are the wear items and their replacement intervals?), and a spare parts list with vendor part numbers. Thinking about Operations \& Maintenance is not an afterthought; it accounts for 50--70\% of a product's lifecycle cost. For the \textbf{Hazards Discussion}: revisit your FMEA with updated RPN scores. Discuss the highest-RPN failure modes you identified, what design changes you made to mitigate them, and whether those changes actually reduced the risk. Also address any new hazards discovered during prototyping and testing that were not in the original FMEA. See MRD Chapter~7 (Operations \& Maintenance) for serviceability principles and the lifecycle cost figure, and Chapter~9 (End of Life) for disposal planning.
\end{faqbox}

\begin{faqbox}[Q: How honest should the ``Lessons Learned'' section be?]
Very. This section is not about making yourself look good; it is about demonstrating that you grew as engineers. The best lessons learned sections candidly discuss what went wrong, why it went wrong, and what the team would do differently. Instructors can tell the difference between generic filler and genuine reflection. Compare these two examples:

\begin{center}
\renewcommand{\arraystretch}{1.3}
\begin{tabularx}{0.95\textwidth}{>{\bfseries\color{warnred}}c >{\raggedright\arraybackslash}X}
\toprule
\textbf{Quality} & \textbf{Example} \\
\midrule
\textcolor{warnred}{Generic} & ``We learned that communication is important and that we should start earlier.'' \\
\textcolor{tipgreen}{Specific} & ``We underestimated the integration complexity between the sensor subsystem and the control logic, which cost us two weeks. Starting in Sprint~4, we adopted a shared engineering notebook with daily entries, which prevented the miscommunications that had plagued Sprints~1--3. Next time, we would prototype the critical interface first.'' \\
\bottomrule
\end{tabularx}
\end{center}

\noindent The generic version could have been written without doing the project. The specific version demonstrates real engineering growth and gives actionable advice to anyone who reads it.
\end{faqbox}

\begin{warnbox}
Each section must start on a new page (page break) so that it aligns correctly in Gradescope. Forgetting this formatting requirement can cause grading issues. Double-check before submitting.
\end{warnbox}

\begin{tipbox}
For the Team Contribution table, have an honest team discussion before filling it out. Agree on rankings together rather than each person filling it out independently and discovering disagreements at submission time.
\end{tipbox}

\begin{faqbox}[Q: How do I know if my FDR package is complete?]
Apply the \textbf{FDR Litmus Test}. Ask six questions: Could a new engineering team who has never seen this project (1)~understand what the system does and why, (2)~operate it using only your Quick-Start Guide, (3)~identify what works and what does not from the VCRM, (4)~reproduce any test result from your test procedures, (5)~replace a failed component using the maintenance plan and spare parts list, and (6)~continue improving the design using your lessons learned and future work recommendations? Each ``no'' represents a gap in your documentation. FDR answers ``\textit{How well did} it work?'' It demands honest, evidence-based accounting, not a sales pitch.
\end{faqbox}

\begin{protip}
\textbf{Prepare an as-designed vs.\ as-built comparison.} Your final report should include side-by-side comparisons of your CDR (as-designed) BOM vs.\ actual (as-built) BOM, and your estimated budget vs.\ actual expenditures. Every component substitution, addition, or deletion should be documented with a rationale. If you saved a CDR baseline snapshot during Sprint~3 (see the CDR tips above), this comparison is easy. If you did not, this will be a painful reconstruction.
\end{protip}

\begin{tipbox}
\textbf{Include O\&M materials in the appendix.} Your final report should include three operations and maintenance deliverables: (1)~a one-page \textbf{Quick-Start Guide} with numbered steps and photos that enables your instructor to set up and operate the system in 10 minutes, (2)~a \textbf{maintenance plan} identifying wear items, replacement intervals, and a spare parts list with part numbers and vendors, and (3)~a \textbf{troubleshooting guide} organized as a symptom-based decision tree for the most common failure modes you discovered. These are not busywork; they demonstrate that you thought about the system beyond your own workbench.
\end{tipbox}

\begin{protip}
\textbf{Trace your results back to the course learning outcomes.} The final report is your opportunity to demonstrate growth across all 10 CLOs listed in the syllabus. As you write each section, ask: which CLO does this evidence support? For example, your design evolution narrative demonstrates CLO~1 (apply the engineering design process) and CLO~4 (build a functional prototype with rapid iterations); your testing summary demonstrates CLO~7 (design experiments) and CLO~9 (collect and analyze data). Making these connections, even informally, strengthens your report and helps your instructor see the full scope of your learning.
\end{protip}

%=============================================================================
